%% JetBrains-tailored CV — Senior Software Engineer, AppGlass (JetBrains Innovation Hub)
%% Based on template_ing.tex (accurate data); platform material restored from template.tex / data/cv.ts
%% Target: https://job-boards.eu.greenhouse.io/jetbrains/jobs/4830653101
\documentclass[10pt,a4paper,sans]{moderncv}
% Load icons before the style so the classic style can measure a bullet symbol
% that exists in the free fontawesome5 set (default \faCircle[regular] is Pro-only)
\pdfmapfile{+fontawesome5.map}
\RequirePackage{fontawesome5}
\renewcommand*{\listitemsymbol}{\faAngleRight}
\moderncvstyle{classic}
\moderncvcolor{orange}
% \faGlobeAmericas is Pro-only; use the free solid globe
\renewcommand*{\homepagesymbol}{{\color{homepage}\small\faGlobe}~}

\usepackage[scale=0.80]{geometry}
\setlength{\hintscolumnwidth}{3cm}
\microtypesetup{expansion=false} % fa5 Type1 fonts are not expandable; keep kerning/protrusion

\name{Giona}{Granchelli}
\title{Senior JVM \& Platform Engineer | Developer Tooling, Distributed Systems \& Agentic AI}
\phone[mobile]{+31~(0)6~427433~16}
\email{giona.granchelli@gmail.com}
\social[linkedin]{giona-granchelli}
\social[github]{GionaGranchelli}
\homepage{gionag.com}
\photo[56pt][0.4pt]{face.jpg}
\quote{Hands-on JVM engineer with 10+ years designing, building, and operating distributed systems in banking, fintech, and cloud-native environments. Deep Java/Kotlin engineering across Spring Boot, Kafka, CQRS, and Kubernetes-based platforms, paired with developer tooling and governed AI-assisted engineering workflows. Currently Chapter Lead Java \& VueJS at ABN AMRO: technically hands-on while leading approximately 45 engineers across 5 teams in a regulated banking environment.}
%----------------------------------------------------------------------------------
% content
%----------------------------------------------------------------------------------
\begin{document}
\vspace*{-2cm}
\makecvtitle

\section{Core Skills}
\cvitem{JVM \& Backend Engineering}{
Java, Kotlin, JVM, Spring Boot, REST APIs, microservices, gRPC, enterprise backends
}
\cvitem{Developer Tooling}{
IntelliJ Platform plugins, agentic coding workflows, CI/CD, developer productivity
}
\cvitem{Distributed Systems}{
Kafka, CQRS, event-driven architecture, Axon Framework, data pipelines
}
\cvitem{Cloud \& Platform}{
Kubernetes, Docker, Azure AKS, AWS, cloud-native migration, deployment patterns
}
\cvitem{Reliability}{
Observability, incident investigation, root-cause analysis, production reliability
}
\cvitem{AI Agents}{
Governed agentic workflows, approval controls, auditable execution
}

\section{Professional Experience}
\cventry{Oct 2024 - Now}{Chapter Lead Java \& VueJS | Software Engineer IV}
{ABN AMRO - BCDB}
{Amsterdam, Netherlands}
{}{\textbf{Copilot Champion \& AI Adoption Lead.} Technically hands-on JVM engineer who also operates at significant organizational scope: Chapter Lead for approximately 45 engineers across 5 teams, accountable for engineering capability, technical standards, cross-team delivery, hiring, and performance management in a regulated banking environment.
\begin{itemize}
 \item Hands-on engineering across the department's JVM estate: reconciliation jobs, ETL pipelines, and legacy Java, Struts, and Spring Boot services, plus the developer tooling that supports them.
 \item Migrated 10 legacy applications (old Java, Struts, Spring Boot) using an agentic migration workflow, cutting an estimated 18+ months of effort to approximately 4 months.
 \item Established a shared repository of reusable agentic workflows for backend, frontend, and infrastructure teams, bringing governed AI-assisted engineering into everyday delivery.
 \item Built agentic workflows integrating Azure DevOps for incident investigation and management, and an end-to-end software-delivery workflow that uses requirements and work-item context to generate and refine stories, implement changes, open PRs, perform reviews, and update engineering systems.
 \item Built the governance framework underpinning all agentic workflows: approval controls, auditable execution, and governed workflows with human-in-the-loop review.
 \item Led modernization initiatives for legacy services, defining migration paths toward cloud-native architectures on Azure AKS, and designed Kubernetes-based platform patterns, CI/CD pipelines, and developer tooling to standardize containerization and deployment across teams.
 \item Built and improved reconciliation jobs validating financial data consistency across multiple banking systems, and improved reliability, observability, and operational robustness of batch and integration flows handling critical banking workloads.
 \item Acted as a technical multiplier for multiple teams by introducing reusable engineering patterns, paved-road style tooling, and migration guidance.
 \item Stack: Kubernetes, Azure, AKS, Azure DevOps, Java, TypeScript, Kafka, agentic AI workflows
\end{itemize}
}

\cventry{Sep 2020 - Sep 2024}{Full-stack | Software Engineer IV}
{ABN AMRO - Strategy and Innovation}
{Amsterdam, Netherlands}
{}{\textbf{PayDay:} core engineer for a fintech platform enabling instant payouts for gig-economy workers and freelancers.
\begin{itemize}
 \item Designed and implemented Kotlin and Spring Boot backend services for payment and payout workflows in a regulated financial environment.
 \item Contributed to architecture decisions around transaction handling, system integration, and scalable service boundaries.
 \item Stack: Kotlin, Spring Boot, Vue.js, Android, Docker, Kubernetes, AWS S3, PostgreSQL, GitLab CI
\end{itemize}
}

\cventry{May 2020 - Jul 2020}{Software Engineer}
{Ximedes}
{Haarlem, Netherlands}
{}{Project for the Dutch public transportation sector (NS): end-to-end solution enabling travel with QR codes instead of conventional tickets.
\begin{itemize}
 \item Built backend and integration components for ticketing workflows; contributed real-time features using WebSockets.
 \item Stack: Kotlin, Ktor, REST, WebSocket, ReactJS, TypeScript, Redis
\end{itemize}
}

\cventry{Oct 2019 - Apr 2020}{Full Stack Developer}
{BLOX - BTC Direct}
{Amsterdam / Nijmegen, Netherlands}
{}{Cryptocurrency trading platform serving approximately 500k daily users.
\begin{itemize}
 \item Maintained and extended microservices for trading, account management, and analytics on a distributed event-driven architecture (Axon Framework, CQRS).
 \item Improved stability and performance across multiple backend services.
 \item Stack: Java 8, Kotlin, Spring Boot, REST, gRPC, Axon Framework, Node.js, ReactJS, MariaDB, Docker, Kubernetes, GitLab CI
\end{itemize}
}

\cventry{Jan 2019 - Sep 2019}{Senior Software Developer}
{WoodWing B.V.}
{Zaandam, Netherlands}
{}{Elvis, a scalable distributed Digital Asset Management system for creating, managing, and distributing digital assets.
\begin{itemize}
 \item Implemented an SSO solution integrating third-party and internal applications (AWS Cognito, Okta) and helped split a monolith by extracting auth capabilities.
 \item Worked with a custom CQRS-style approach built around a modified use of Elasticsearch.
 \item Stack: Java 8, Spring Framework, Elasticsearch, AWS, Jenkins, Docker, REST, AngularJS, Node.js, Okta
\end{itemize}
}

\cventry{Jul 2017 - Jan 2019}{Full Stack Developer}
{Trifork B.V.}
{Amsterdam, Netherlands}
{}{Multiple projects in small agile consultancy teams, combining backend delivery, frontend work, and rapid onboarding into new domains.
\begin{itemize}
 \item Contributed to IBE, a cloud-based teacher-student exam platform used in Switzerland; migrated and rebuilt CI/CD pipelines from Jenkins to GitLab.
 \item Stack: J2EE, Kotlin, Spring Boot, REST, gRPC, Axon Framework, Node.js, ReactJS, Docker, Kubernetes, Jenkins, GitLab CI
\end{itemize}
}

\cventry{Oct 2016 - Mar 2017}{Research Scholarship: Tactics (EU project)}
{University of L'Aquila}
{L'Aquila, Italy}
{}{European project Tactics with Thales Chieti and Finmeccanica (Leonardo): SOAP-based services over a new ZigBee military tactical network for NATO troops.
\begin{itemize}
 \item Stack: J2EE, Spring Framework, Spring Boot, SOAP, REST, Hibernate, JQuery, DDS, WebDDS
\end{itemize}
}

\cventry{Mar 2016 - Sep 2016}{Research Scholarship: USRA}
{University of L'Aquila}
{L'Aquila, Italy}
{}{Enterprise system for the post-earthquake rebuilding of L'Aquila with USRA (Ufficio Speciale Ricostruzione dell'Aquila): inspection, review, request, classification, and refund of real estate. Developed from scratch.
\begin{itemize}
 \item Stack: J2EE, Hibernate, Spring Framework, Maven, SVN, REST, JQuery, MySQL, HTML5, CSS3
\end{itemize}
}

\section{Selected Projects}
\cvitem{IntelliAiBridge (IntelliJ Platform plugin, JetBrains Marketplace: \href{https://share.google/5Kz9f7VJP58gyIlbQ}{AiBridge})}{
IntelliJ Platform plugin/integration enabling approved external agentic tooling to use enterprise Copilot capabilities while preserving enterprise authentication, authorization, and network/data boundaries. External agents talk to the plugin through the standard OpenAI protocol, which routes requests through the authenticated Copilot session already running inside the IDE.
}
\cvitem{TramAI (open source, Apache-2.0, \href{https://tramai.dev}{tramai.dev})}{
A JVM runtime for governed AI applications: portable model providers, multi-agent workflows, approval controls, and auditable execution. Designed for production systems where portability, governance, and operational visibility are requirements rather than optional extras (\href{https://github.com/GionaGranchelli/tramAI}{github.com/GionaGranchelli/tramAI}).
}
\cvitem{QMD / OpenClaw}{
Kotlin Multiplatform/Ktor control surface for multi-model agent orchestration; QMD provides surgical local retrieval over thousands of documents.
}

\section{Talks \& Teaching}
\cvitem{June 2026}{Guest lecturer, Vrije Universiteit Amsterdam (VU Amsterdam): industry lecture on AI-assisted and agentic software engineering.}
\cvitem{2026}{Speaker, Ktconf 2026 (Kotlin conference).}

\section{Education}
\cventry{Sep 2015 - Mar 2017}{Master Degree in Computer Science}
{University of L'Aquila}
{L'Aquila (AQ), Italy}
{maximum score, cum laude}{
Thesis: Architecture Recovery of Microservice-based Systems.
}

\cventry{Dec 2013 - Jun 2014}{First Level Postgraduate Master Degree in Web Technologies}
{University of L'Aquila}
{L'Aquila (AQ), Italy}
{}{}

\cventry{Sep 2009 - Jul 2013}{Bachelor Degree in Computer Science}
{University of L'Aquila}
{L'Aquila (AQ), Italy}
{99/110}{}

\section{Publications}
\cvitem{[1]}{MicroART: A Software Architecture Recovery Tool for Maintaining Microservice-based Systems. In Proceedings of the 14th International Conference on Software Architecture (ICSA). IEEE, 2017.}
\cvitem{[2]}{Towards Recovering the Software Architecture of Microservice-based Systems. In IEEE International Workshop on Architecting with MicroServices (AMS), April 2017.}

\end{document}
